\section{Ingestion des données dans HDFS}

L'arborescence \texttt{/projet\_ia} sépare trois zones : \texttt{raw}
(donnée source, jamais modifiée), \texttt{processed} (les 3 tables
\texttt{offres}/\texttt{metiers}/\texttt{competences} construites par
Spark, voir la section Modèle offres-métiers-compétences plus loin) et
\texttt{output}
(exports agrégés pour Zeppelin/DAS).

\begin{lstlisting}
/projet_ia/
├── raw/                                    (CSV source)
├── processed/
│   ├── offres/
│   ├── metiers/
│   └── competences/
└── output/
\end{lstlisting}

Créée et peuplée via deux scripts versionnés
(\texttt{infra/scripts/hdfs-bootstrap.sh} et
\texttt{infra/scripts/hdfs-verify.sh}), pas de commande ad-hoc non
tracée :

\begin{lstlisting}
$ hdfs dfs -mkdir -p /projet_ia/raw /projet_ia/processed/offres \
    /projet_ia/processed/metiers /projet_ia/processed/competences \
    /projet_ia/output
$ hdfs dfs -put -f ai_jobs_market_2025_2026.csv /projet_ia/raw/
\end{lstlisting}

\textbf{Contrôles effectués} :

\begin{lstlisting}
$ hdfs dfs -ls -R /projet_ia
-rw-r--r--   2 hadoop supergroup     383900  .../raw/ai_jobs_market_2025_2026.csv

$ hdfs dfs -du -h /projet_ia
374.9 K  749.8 K  /projet_ia/raw

$ hdfs dfs -cat /projet_ia/raw/ai_jobs_market_2025_2026.csv | wc -l
1501
\end{lstlisting}

1501 lignes = 1500 lignes de données + 1 ligne d'en-tête, cohérent avec
\texttt{docs/dataset-schema.md} (1500 offres). Le facteur 2 sur
\texttt{-du} (374.9 K de contenu réel pour 749.8 K de "disk space
consumed") s'explique par \texttt{dfs.replication=2} configuré sur le
cluster (\texttt{infra/manifests/10-hdfs/configmap.yaml}) : chaque bloc
est stocké deux fois, sur les deux DataNodes.

La zone \texttt{raw} n'est jamais modifiée après dépôt : toute
transformation (nettoyage, typage, construction des tables) passe par
Spark et écrit exclusivement dans \texttt{processed}, jamais en place —
ce qui permet de rejouer le pipeline à l'identique depuis la source si
une règle de nettoyage doit être révisée.
